\chapter{Introducción}
\label{chap:introduccion}

\section*{Pregunta de apertura}

¿Puede un sistema aprender a ordenar acciones fuera de muestra y convertir después ese orden en una
cartera útil? La pregunta parece única, pero contiene dos problemas. Ordenar exige que las acciones
con mayor puntuación obtengan, en promedio, mejores retornos futuros que las de menor puntuación.
Invertir exige además elegir cuántas comprar, cuánto asignar, cuándo vender y cuánto efectivo
mantener. Una señal puede resolver el primer problema y fracasar en el segundo.

Esa separación organiza todo el trabajo. El modelo se selecciona únicamente por su capacidad de
ordenación transversal. Las reglas de cartera se estudian después, con la señal congelada. Así, una
mejora económica no puede reescribirse retrospectivamente como mejora predictiva y la era
2025--2026 permanece reservada desde el inicio.

\section{Problema y relevancia}

La selección cuantitativa de acciones combina tres dificultades. Primero, la información disponible
cambia con el tiempo: usar hoy la composición actual de un índice o un dato fundamental antes de su
publicación produce una ventaja inexistente. Segundo, probar muchas configuraciones aumenta la
probabilidad de elegir ruido. Tercero, el resultado observado depende de una cadena de decisiones
posteriores al modelo. Evaluar solo la rentabilidad final no permite saber en qué eslabón apareció la
mejora o el fallo.

El sistema estudiado afronta esas dificultades con un panel \textit{point-in-time}, cohortes
cerradas y una única ruta experimental. Cinco agentes ---\texttt{quality}, \texttt{value},
\texttt{growth}, \texttt{momentum} y \texttt{risk}--- producen rankings especializados; un
meta-agente aprende a combinarlos con información que ya estaba cerrada en cada fecha. La selección
secuencial modifica una variable cada vez y conserva únicamente cambios respaldados por Rank-IC
robusto. Después se calculan atribución, robustez, carteras y perfiles como evidencia informativa.

\section{Dos objetivos}

\textbf{Objetivo 1.} Determinar si el sistema aprende una ordenación transversal con capacidad
predictiva fuera de muestra y si esa capacidad resiste cambios de régimen y controles contra azar,
fuga temporal y factores conocidos. Este objetivo no presupone que el meta-agente deba superar a
todos los especialistas: esa comparación forma parte del resultado.

\textbf{Objetivo 2.} Medir cuánto alteran las reglas de cartera la traducción económica de una señal
ya congelada. Este objetivo se divide al interpretar, no al ejecutar: (2a) comprobar si las reglas
generan diferencias materiales dentro de la ventana de selección; y (2b) observar si la cartera
elegida conserva el resultado en la era reservada. La primera afirmación puede sostenerse con toda la
rejilla; la segunda queda limitada por el reducido número de cohortes reservadas.

Los cinco criterios que permiten juzgar ambos objetivos son: integridad temporal de datos y
etiquetas; mejora predictiva medible y comparable; robustez frente a explicaciones nulas;
separación entre señal y cartera; y límites explícitos sobre cobertura, multiplicidad y capacidad.
No son nuevas preguntas de investigación, sino condiciones que una respuesta defendible debe
satisfacer.

\section{Aportación}

La aportación no consiste en presentar una cartera como recomendación de inversión. Consiste en un
procedimiento auditable para separar aprendizaje predictivo e implementación económica, junto con
un caso empírico completo. El trabajo muestra qué aprendió el sistema, qué especialista concentra la
información, qué contrastes supera, qué reglas de cartera mueven el resultado y qué cambia cuando la
misma señal se evalúa en una ventana que no intervino en ninguna decisión.

La distinción numérica también es importante. El sistema emplea \NumPredictores{} predictores
declarados; el catálogo contiene \NumParametrosCatalogo{} parámetros científicos y de cartera; y la
auditoría de variables produce \NumDiagnosticos{} diagnósticos porque incluye columnas derivadas.
Son tres objetos distintos, no tres versiones del mismo recuento.

\section{Recorrido del documento}

El capítulo siguiente sitúa la pregunta en la literatura. Después se presentan datos y agentes, se
formaliza el protocolo, se evalúa la señal y se estudia su traducción a cartera. Limitaciones y
conclusiones cierran el argumento; los anexos conservan únicamente lo necesario para reproducir el
estudio y verificar el catálogo cerrado.

\section*{Respuesta y transición}

El problema de investigación queda planteado como una cadena causal, no como una única cifra de
rentabilidad. Para que esa cadena sea creíble, el siguiente capítulo delimita qué aportan Rank-IC,
la validación temporal y la teoría de construcción de carteras, y qué hueco concreto cubre este
trabajo.
